% Round 18. U2 branch. Candidate: requires a bound independent review.
\section{Ramified base change of positive fourfold models}
\label{r18rb:section}

This section treats new cycles over a ramified field by constructing a
new reduced strict semistable witness.  The calculation is algebraic on
its actual strata.  It uses neither an assertion that those cycles
already descend nor an unspecified proper pushforward of actual forms.
The geometric input about toroidal factorization is the corrected
\emph{weak} factorization theorem, not the withdrawn strong version.

\subsection{The two precise toroidal inputs}

We use the following published inputs in their indicated scopes.
\begin{enumerate}
\item Takeshi Saito, \emph{Weight spectral sequences and independence
of $\ell$}, Lemmas 1.2.4--1.2.5, author manuscript
\url{https://www.ms.u-tokyo.ac.jp/~t-saito/pp/wmr2r.pdf},
printed pp.~15--16.  If $\mathcal X/R$ is strict semistable, its
components have ideals $I_i=\mathcal O_{\mathcal X}(-D_i)$,
and $R_L/R$ has ramification index $e$ and uniformizer $t$, then the
\emph{actual blowup} of $\mathcal X_{R_L}$ in the coherent ideal
\[
 \mathfrak a_e=
 \prod_{0\le j<k\le m}\ \prod_{h=1}^{e-1}
       \bigl(I_{j+1}\cdots I_k\mathcal O_{\mathcal X_{R_L}}
                      +t^h\mathcal O_{\mathcal X_{R_L}}\bigr)
 \tag{RB18.1}\label{r18rb:saito-ideal}
\]
is reduced strict semistable.  For $e=1$ the ideal is the unit ideal.
The proof is by the displayed monoid charts and integral unimodular
subdivision.  It imposes no tame-ramification assumption.

\item The integral-conical-complex weak factorization theorem:
a proper birational toroidal map between nonsingular strict toroidal
embeddings can be factored into smooth toroidal blowups and blowdowns;
if the endpoints are projective, the intermediate varieties can be
chosen projective.  Here we use the corrected conical-complex statement
of Kenji Matsuki, \emph{Erratum to ``A note on the factorization theorem
of toric birational maps after Morelli and its toroidal extension''},
item~2, \url{https://arxiv.org/pdf/math/0006038}, together with the
projectivity construction in Abramovich--Karu--Matsuki--W\l odarczyk,
\emph{Torification and factorization of birational maps},
Theorem~2.7.1 and its proof, JAMS \textbf{15} (2002), 551--552,
\url{https://personal.math.ubc.ca/~karu/papers/semist.pdf}.
\end{enumerate}

We spell out the characteristic issue in the second input.  We do
\emph{not} invoke the general characteristic-zero weak factorization
theorem for an arbitrary birational map of special-fiber varieties.
The input used here is its already toroidal, purely integral
conical-complex part: the factorization consists of smooth star
subdivisions and inverse operations in the integral cone complexes.
These combinatorial operations do not depend on a ground-field
characteristic.  A smooth star subdivision is realized, over any field,
by blowing up the corresponding smooth closed stratum; its monomial
charts and their gluing are the same integral charts in every
characteristic.  A rational strictly convex support function realizes
the asserted projectivity.  A common toroidal graph and its projective
monomial resolution replace elimination of indeterminacy at the start
of the projectivity argument.  No resolution of arbitrary
positive-characteristic singularities enters that argument.

In our application every boundary component and the lattices of the
normal torus characters are individually labelled.  After a finite
unramified extension, if needed, all connected old strata are
geometrically integral.  Thus these are strict conical complexes
without self-identifications; each smooth-star center is a native
smooth closed stratum, or its disjoint connected components.  Applying
the integral algorithm and then its monomial realization gives a
specified finite chain of smooth projective blowups and blowdowns
on the varieties under consideration.  This explanation does not
assert arbitrary-field weak factorization for nontoroidal maps.

\subsection{Every new stratum has an actual toroidal comparison}

Let $\mathcal X/R$ be a regular projective reduced strict semistable
model of relative dimension $n$.  Write $D_I=\bigcap_{i\in I}D_i$.
Fix a finite extension $L/K$ and let $\mathcal Y/R_L$ be the blowup
\eqref{r18rb:saito-ideal}.  Its generic fiber is exactly $X_L$.
We first work after a residue extension splitting all connected old
strata, so that each of their components is geometrically integral.
The field descent at the end will retain its full degree.

\begin{lemma}[Toroidal comparison for each actual closed stratum]
\label{r18rb:strata}
Let $S$ be a geometrically integral component of an actual closed
stratum formed by $q$ distinct components of $\mathcal Y_k$.
Then there is an old geometrically integral closed stratum $D_I$,
with $s=|I|\ge q$, such that $S\to\mathcal X_k$ factors through
$D_I$ and dominates it.  Put $b=s-q$.  There are line bundles
$\mathcal L_1,\ldots,\mathcal L_b$ on $D_I$ for which $S$ and
\[
 \mathcal P_S=
   \prod_{\nu=1}^{b}{}_{D_I}
       \mathbf P_{D_I}(\mathcal O\oplus\mathcal L_\nu^\vee)
 \tag{RB18.2}\label{r18rb:bundle}
\]
are connected by a finite chain of smooth projective blowups and
blowdowns.  The projective bundles use the quotient convention;
for $b=0$ the right side means $D_I$.
The birational comparison and the factorization are toroidal and
are identities on the same normal torus torsor over
\[
 D_I^\circ=D_I\setminus\bigcup_{j\notin I}D_j.
\]
No birational factorization theorem for arbitrary varieties is used.
\end{lemma}

\begin{proof}
\emph{Step 1: the integral orbit lattice.}
At a point of an old $s$-fold stratum, strict semistability gives an
actual etale chart $x_1\cdots x_s=\pi$, with extra smooth coordinates.
After base change write $\pi=u t^e$ and absorb the unit $u$ in one
coordinate.  The local monomial character lattice and its dual are
\[
 M_I=\frac{\mathbf Z^{I}\oplus\mathbf Z[t]}
               {\mathbf Z((1)_{i\in I},-e)},\qquad
 N_I=\{((a_i)_{i\in I},h):\sum_i a_i=eh\}.
 \tag{RB18.3}\label{r18rb:lattices}
\]
The Saito subdivision has height-one rays
$((a_i),1)$, with $a_i\ge0$ and $\sum_i a_i=e$.
Every cone is unimodular in $N_I$.
An actual $q$-fold stratum corresponds to a $q$-dimensional cone
$\sigma$; the union of the supports of its rays determines $I$.
Its relative interior maps to the relative interior of the old
cone for $D_I$, so $S$ dominates $D_I$.
The character lattice of the dense orbit in the fiber is
\[
                  M_\sigma=M_I\cap\sigma^\perp,
                       \qquad\operatorname{rank}M_\sigma=s-q=b.
 \tag{RB18.4}\label{r18rb:orbit-lattice}
\]
Unimodularity makes the orbit a split torus of this rank, with no
hidden finite cover of the old stratum.  Equivalently, the local
monoid ring of the open orbit is the group ring of this free
lattice over the coordinate ring of $D_I^\circ$.
All these descriptions are compatible on faces of the old
cone complex, hence describe the actual global stratum, not a
formal substitute for it.  A connected component of an old
stratum is treated separately if the intersection is disconnected.

\emph{Step 2: the torsor and the extending line bundles.}
Choose a lattice basis $\chi_1,\ldots,\chi_b$ of $M_\sigma$ and
fixed integral representatives
$((c_{i\nu})_{i\in I},d_\nu)$ in the numerator of
\eqref{r18rb:lattices}.  On the local open orbit the corresponding
coordinate is the residue of the valuation-zero monomial
\[
                           t^{d_\nu}\prod_{i\in I}x_i^{c_{i\nu}}.
 \tag{RB18.5}\label{r18rb:characters}
\]
Negative exponents here are permitted: these are invertible
characters on the open orbit and rational functions on its
compactification, not purported global regular functions.

If a local defining equation of $D_i$ changes from $x_i$ to
$g_i x_i$, the coordinate \eqref{r18rb:characters} changes by
$\prod_i(g_i|_{D_I})^{c_{i\nu}}$.  Therefore the torus torsor
extends from $D_I^\circ$ to the product of the nonzero line-bundle
spaces for
\[
                  \mathcal L_\nu=
                  \bigotimes_{i\in I}
                  \bigl(\mathcal O_{\mathcal X}(D_i)|_{D_I}\bigr)
                                      ^{\otimes c_{i\nu}}.
 \tag{RB18.6}\label{r18rb:line-bundles}
\]
With the quotient convention,
$\mathbf P(\mathcal O\oplus\mathcal L_\nu^\vee)$ contains
$\mathcal L_\nu^\times$ as the open set where its two quotient
coordinates are nonzero.  This verifies the convention in
\eqref{r18rb:bundle}.  A different representative of a character
changes its extension by a boundary line bundle; a fixed choice
has been made, and both extensions identify the same torus torsor
on $D_I^\circ$.  No root of a transition unit is extracted, even
when $e$ is divisible by the residue characteristic.

The local monoid charts show that this torsor is precisely the
open subset of $S$ avoiding the other new components.  Indeed a
new boundary divisor is an adjacent toric orbit divisor, and the
preimage of any old extra component is a union of such divisors.
The transitions just computed depend only on their restrictions
to $D_I$, since the map from $S$ factors through $D_I$.  Thus the
identification is an actual algebraic identification on the common
open subset; it does not rely on choosing a formal tubular
neighborhood of $D_I$.

\emph{Step 3: toroidal compactifications and their factorization.}
The complement of the common open in $S$ is the restriction of the
other components of the reduced SNC fiber of $\mathcal Y$.
It is a strict normal-crossings divisor.  On $\mathcal P_S$ the
complement is the union of the zero and infinity sections of its
line-bundle factors and the inverse image of the old boundary of
$D_I$.  This too is a strict normal-crossings divisor.
The induced birational map between these two proper smooth
compactifications is toroidal: on the actual charts its fiber
coordinates are the integral monomials
\eqref{r18rb:characters}, while its old boundary coordinates and
smooth base coordinates are pulled back from $D_I$.
In particular there is one fixed common open torus torsor; merely
being locally toric with unrelated choices of open tori would
not suffice.

Apply the corrected integral-conical-complex weak factorization
input above to this toroidal birational map.  Smooth stars are
realized by blowups of the indicated native smooth closed strata,
and the projectivity refinement keeps all intermediate varieties
projective.  This gives the chain asserted in the lemma.  The
construction does not claim that all blowups precede all blowdowns.
\end{proof}

\subsection{Fourfold numerical positivity and the double-curve gate}

Call a smooth projective fourfold \emph{universally middle-positive}
if, after every finite extension of its native finite field and on
every resulting connected component, its native middle numerical
pairing has negative index $\rho-1$, where $\rho=\dim N^1$.
By S18 this is equivalent to the numerical primitive middle
Hodge positivity for every ample class on that component.
All other required numerical Lefschetz and Hodge-index assertions
on a fourfold follow from divisor Hodge index as in S18.

Projective bundles of total dimension four over smooth projective
bases of dimension at most three are universally middle-positive,
by S18.  The same holds for their specified smooth blowups and
blowdowns.  Smooth four-dimensional quadrics are also universally
middle-positive: the frozen C217 quadric middle-form lemma in every characteristic
(with $a=2$, extending QF17) gives positive primitive form, and the ample square gives the
remaining positive line.  Equivalently, after splitting, the two
middle plane classes have intersection matrix $\operatorname{diag}(1,1)$;
$N^1$ has rank one.  Pullback to a splitting field and the full native
degree pairing retain this positivity on the native subspace.
S18's computed exceptional blocks then give the same assertion for
specified smooth blowups and blowdowns of these quadric fourfolds.

\begin{theorem}[Arbitrary finite base change of a positive fourfold witness]
\label{r18rb:fourfold}
Let $X/K$ be a smooth projective geometrically integral fourfold
with a native regular projective reduced strict semistable model
$\mathcal X/R$.  Assume:
\begin{enumerate}
\item every component $D_i$ is universally middle-positive in the
precise numerical sense just defined;
\item on every double stratum, over every finite residue extension,
every numerically zero rational curve has zero geometric
$\ell$-adic cycle class.
\end{enumerate}
Then, for every finite valued extension $L/K$ and every specified
geometrically homologically trivial rational cycle on $X_L$, strong
G holds with exact original support and all ordinary actual compact
tests over $L$.  In particular the fourfold conclusion of HP18 is
stable on the full cycle groups under arbitrary finite extension.
\end{theorem}

\begin{proof}
Fix $L$ and a cycle on $X_L$; no descent of this cycle is assumed.
Choose a further finite unramified extension $L'/L$ splitting the
finitely many connected components of all old strata if necessary.
This does not change the ramification index in Saito's construction.
The hypotheses explicitly persist under its residue extension.
Let $\mathcal Y/R_{L'}$ be the actual projective reduced strict
semistable model \eqref{r18rb:saito-ideal}.

A new component has $q=1$ in Lemma~\ref{r18rb:strata}.
If $s=1$, its toroidal comparison is a chain of smooth projective
blowups and blowdowns from the corresponding old fourfold component.
If $s\ge2$, the comparison bundle has positive torus rank $s-1$;
its last factor is a $\mathbf P^1$-bundle over a smooth projective
threefold.  It is therefore middle-positive by S18.  In both cases
S18's exceptional-pairing computation transports middle positivity
through the entire chain.  Every new fourfold component consequently
satisfies numerical middle Hodge positivity for the restriction of
one actual ample divisor on $\mathcal Y$.  All its other nonempty
strata have dimension at most three and satisfy the numerical
Lefschetz and primitive Hodge-index assertions by divisor Hodge
index.  The frozen R6 numerical BGS criterion therefore gives the
numerical local exactness $U_2(\mathcal Y)=0$.

A new double has $q=2$.  If $s=2$, its comparison bundle is the old
threefold double itself, which has the assumed curve comparison.
If $s\ge3$, the bundle has $s-2\ge1$ projective-line factors and
total dimension three.  Its last factor is a projective-line bundle
over a smooth projective surface, hence has curve comparison by
CD17.  That comparison is preserved through all the supplied smooth
blowups and blowdowns, again by CD17.  Thus every actual new double
has the curve-comparison property, including after any further
finite residue extension.  CL17 gives
$K_{\mathrm{comp},\ell}^2(\mathcal Y)=\ker\rho_2$.
DQ17 and ER17 now give the zero refined obstruction for every
specified homologically trivial codimension-two cycle on $X_{L'}$.
The frozen actual reconstruction supplies its exact-support Green
form against every compact actual test.  The other codimensions
are the frozen divisor, one-cycle, and zero-cycle cases.

Apply the just established conclusion on $\mathcal Y$ to the
pullback of the arbitrary specified cycle fixed over $L$.  The frozen finite actual trace
descends with division by the full $[L':L]$, including its unramified
constant-field degree, and with the absolute value extending the
original one on $L$.  This returns the form on the original support
complement over $L$ with all original actual tests.  Since the cycle
over $L$ was arbitrary at the start, the entire asserted cycle
domain over every $L$ is covered.
\end{proof}

\begin{corollary}[Hyperplanes and linear--quadric arrangements]
\label{r18rb:families}
The full fourfold G theorem of HP18 remains true for every new cycle
after every finite valued extension.  The same conclusion holds for
the fourfold NF18 models whose factor degrees all lie in $\{1,2\}$.
Consequently the fourfold global HP18/GF18 height families, and the
fourfold NG18 families from these linear--quadric profiles, have the
stated complete complementary-degree height pairings after
\emph{every} number-field extension.  The restriction ``unramified
at $q$'' is unnecessary for these particular fourfold families.
\end{corollary}

\begin{proof}
The HP18 components are the specified smooth blowups of
$\mathbf P^4$.  The indicated NF18 components are the specified
smooth blowups of smooth hyperplanes or quadric fourfolds.
They are universally middle-positive by the preceding calculation
and S18.  HP18's actual doubles have CD17's curve comparison;
NF18 proves the same property for its complete-intersection
threefold doubles and their specified smooth blowups.  Apply
Theorem~\ref{r18rb:fourfold} anew at every completion above the
formerly distinguished prime $q$.  At all other places the frozen
good or ordinary-quadratic theorem already accepts every finite
extension.  PH16 then supplies the complete height pairings, with
GA17's admissible model lifts, on the full homologically trivial
Chow groups over the new number field.
\end{proof}

\paragraph{Limits and dependencies.}
For nonlinear factor fourfolds of arbitrary degree, the native middle
positivity of the original components is not supplied by NF18; that
paper's direct DG18 proof does not assert it.  In particular the
five-sextic profile is not included in the new extension corollary.
For relative dimension greater than four, this argument does not
supply the required higher numerical positivity of the new
components, and gives no blanket ramified-extension upgrade of
HG18.  General G remains open.  The new conclusion is obtained by
an actual new semistable model, its toroidal stratum comparisons,
S18, the corrected numerical BGS input frozen in R6, CD17, CL17,
DQ17, ER17, and the full-degree finite trace.  The known resolution
and weak factorization theorems are cited as external inputs, and
no priority claim is made for them.
